\documentclass[12pt]{beamer}

% ── Thème et couleurs ─────────────────────────────────────────
\usetheme{default}
\usecolortheme{default}
\setbeamertemplate{navigation symbols}{}
\setbeamertemplate{footline}[frame number]

% Couleurs personnalisées
\definecolor{bleutitre}{RGB}{0,70,140}
\definecolor{bleuaccent}{RGB}{0,102,204}
\definecolor{grisclair}{RGB}{240,245,250}

\setbeamercolor{frametitle}{fg=bleutitre, bg=grisclair}
\setbeamercolor{title}{fg=bleutitre}
\setbeamercolor{subtitle}{fg=bleuaccent}
\setbeamercolor{structure}{fg=bleutitre}
\setbeamercolor{normal text}{fg=black, bg=white}
\setbeamercolor{block title}{fg=white, bg=bleutitre}
\setbeamercolor{block body}{fg=black, bg=grisclair}
\setbeamercolor{alerted text}{fg=bleuaccent}

\setbeamerfont{frametitle}{size=\large, series=\bfseries}
\setbeamerfont{title}{size=\Large, series=\bfseries}

% Trait bleu sous le titre de chaque diapo
\setbeamertemplate{frametitle}{
  \vspace*{0.3em}
  \insertframetitle
  \vspace*{0.2em}
  \textcolor{bleuaccent}{\hrule height 1.2pt}
  \vspace*{0.3em}
}

% ── Packages ──────────────────────────────────────────────────
\usepackage[utf8]{inputenc}
\usepackage[T1]{fontenc}
\usepackage[french]{babel}
\usepackage{amsmath, amssymb}
\usepackage{booktabs}
\usepackage{graphicx}
\usepackage{multirow}
\usepackage{array}
\usepackage{tikz}
\usetikzlibrary{arrows.meta, positioning}

% ── Métadonnées ───────────────────────────────────────────────
\title{Modélisation des comportements de santé\\par apprentissage automatique}
\subtitle{Régression, classification et clustering sur données synthétiques}
\author{Ardacham Mahamat Teguene}
\institute{Master 1 MIASHS ---~Université de Montpellier Paul Valéry\\
\smallskip Tuteurs : Sophie Lèbre \& Jérôme Pasquet}
\date{Juin 2026}

% ══════════════════════════════════════════════════════════════
\begin{document}
% ══════════════════════════════════════════════════════════════

% ── Diapo 1 : Titre ──────────────────────────────────────────
\begin{frame}
\titlepage
\end{frame}

% ── Diapo 2 : Plan ───────────────────────────────────────────
\begin{frame}{Plan de la présentation}
\tableofcontents
\end{frame}

% ══════════════════════════════════════════════════════════════
\section{Contexte et problématique}
% ══════════════════════════════════════════════════════════════

\begin{frame}{Contexte et problématique}

\begin{block}{Question centrale}
\centering
\textbf{Quels facteurs du mode de vie influencent le plus l'état de santé\\
---~ et peut-on les identifier, les modéliser et les expliquer ?}
\end{block}

\bigskip

\textbf{Motivations :}
\begin{itemize}
  \item Les maladies chroniques résultent de comportements quotidiens mesurables
  \item Données de santé de plus en plus disponibles (IoT, dossiers médicaux)
  \item ML peut modéliser des relations complexes là où la statistique classique atteint ses limites
\end{itemize}

\bigskip

\textbf{Seuils médicaux de référence :}
\begin{itemize}
  \item Obésité : IMC $> 30$ (OMS)
  \item Hypercholestérolémie : cholestérol $> 200$ mg/dL (NCEP)
  \item Hypertension : pression systolique $\geq 140$ mmHg (ESC)
\end{itemize}

\end{frame}

% ══════════════════════════════════════════════════════════════
\section{Jeu de données et pipeline}
% ══════════════════════════════════════════════════════════════

\begin{frame}{Le jeu de données ---~\textit{Health \& Lifestyle Dataset}}

\begin{columns}[T]
\begin{column}{0.5\textwidth}
\textbf{Caractéristiques :}
\begin{itemize}
  \item \textbf{100 000 individus}, 15 variables
  \item Dataset \textbf{synthétique} (Kaggle)
  \item Aucune valeur manquante
  \item Aucun outlier (méthode IQR)
\end{itemize}

\bigskip

\textbf{Variables cibles :}
\begin{itemize}
  \item \texttt{cholesterol} $\rightarrow$ régression
  \item \texttt{disease\_risk} $\rightarrow$ classification
\end{itemize}
\end{column}

\begin{column}{0.48\textwidth}
\textbf{Types de variables :}

\smallskip
\begin{tabular}{ll}
\toprule
Type & Variables \\
\midrule
Continues & age, bmi, steps \\
            & sleep, cholesterol \\
            & systolic\_bp, hr \\
Binaires    & smoker, alcohol \\
            & family\_history \\
            & disease\_risk \\
Catégorielle & gender \\
\bottomrule
\end{tabular}
\end{column}
\end{columns}

\end{frame}

% ─────────────────────────────────────────────────────────────
\begin{frame}{Pipeline CRISP-DM}

\begin{center}
\begin{tikzpicture}[
  box/.style={draw=bleutitre, fill=grisclair, rounded corners=4pt,
              minimum width=3.2cm, minimum height=0.7cm,
              font=\small\bfseries, text=bleutitre},
  arr/.style={->, thick, color=bleuaccent}
]
\node[box] (eda)  at (0,0)    {EDA};
\node[box] (fe)   at (3.5,0)  {Feature Engineering};
\node[box] (reg)  at (0,-1.4) {Régression};
\node[box] (cla)  at (3.5,-1.4) {Classification};
\node[box] (clu)  at (0,-2.8) {Clustering};
\node[box] (int)  at (3.5,-2.8) {Interprétabilité};

\draw[arr] (eda) -- (fe);
\draw[arr] (fe)  -- (reg);
\draw[arr] (fe)  -- (cla);
\draw[arr] (reg) -- (clu);
\draw[arr] (cla) -- (int);
\draw[arr] (clu) -- (int);
\end{tikzpicture}
\end{center}

\smallskip
\begin{itemize}
  \item Approche \textbf{itérative} : chaque étape informe la suivante
  \item EDA $\rightarrow$ anticipe les problèmes de signal
  \item Feature engineering $\rightarrow$ 3 variables métier créées, data leakage détecté
\end{itemize}

\end{frame}

% ══════════════════════════════════════════════════════════════
\section{Analyse exploratoire (EDA)}
% ══════════════════════════════════════════════════════════════

\begin{frame}{EDA ---~Résultats clés}

\begin{columns}[T]
\begin{column}{0.48\textwidth}
\begin{block}{Distributions}
\begin{itemize}
  \item Skewness $< |0.02|$ pour toutes les variables continues
  \item Distributions \textbf{quasi-uniformes}
  \item Aucune transformation nécessaire
\end{itemize}
\end{block}

\medskip

\begin{block}{Classe cible}
\begin{itemize}
  \item Déséquilibre \textbf{75\% / 25\%}
  \item Modèle naïf $\rightarrow$ 75\% d'accuracy sans rien apprendre
  \item Métriques adaptées : F1, AUC-ROC
\end{itemize}
\end{block}
\end{column}

\begin{column}{0.48\textwidth}
\begin{block}{\alert{Résultat central}}
\begin{itemize}
  \item Corrélations $< |0.01|$ entre \textbf{toutes} les variables
  \item Variable la plus corrélée à \texttt{disease\_risk} : \texttt{resting\_hr} avec $r = 0.005$
  \item \textbf{Signal linéaire quasi-nul}
\end{itemize}
\end{block}

\medskip

\textbf{Implication :} orientation vers les méthodes non-linéaires (RF, XGBoost)
\end{column}
\end{columns}

\end{frame}

% ══════════════════════════════════════════════════════════════
\section{Modélisation supervisée}
% ══════════════════════════════════════════════════════════════

\begin{frame}{Régression ---~Protocole et modèles}

\textbf{Objectif :} prédire \texttt{cholesterol} (mg/dL)

\medskip

\textbf{Protocole :}
\begin{itemize}
  \item Split 80/20 ---~Validation croisée K-Fold ($k=5$)
  \item Tuning par \texttt{RandomizedSearchCV} (\texttt{n\_iter=20})
  \item Métriques : $R^2$, RMSE, MAE
\end{itemize}

\medskip

\textbf{6 modèles ---~logique de spectre de complexité :}

\smallskip
\begin{tabular}{lll}
\toprule
Modèle & Type & Spécificité \\
\midrule
Régression linéaire & Baseline & Référence \\
Ridge ($L_2$) & Régularisée & Réduit les coefficients \\
Lasso ($L_1$) & Régularisée & Annule des coefficients \\
ElasticNet & Régularisée & $L_1 + L_2$ \\
Random Forest & Ensembliste & Non-linéaire, parallèle \\
XGBoost & Ensembliste & Non-linéaire, séquentiel \\
\bottomrule
\end{tabular}

\end{frame}

% ─────────────────────────────────────────────────────────────
\begin{frame}{Régression ---~Résultats}

\begin{center}
\small
\begin{tabular}{lrrrr}
\toprule
Modèle & $R^2$ & RMSE & MAE & CV-$R^2$ \\
\midrule
Linéaire (baseline) & 0.0001 & 43.33 & 37.59 & $\approx 0$ \\
Ridge (tuné)        & 0.0001 & 43.33 & 37.59 & $\approx 0$ \\
Lasso (tuné)        & $-0.000$ & 43.33 & 37.60 & $\approx 0$ \\
ElasticNet          & $-0.000$ & 43.33 & 37.60 & $\approx 0$ \\
Random Forest       & $-0.005$ & 43.44 & 37.67 & $\approx 0$ \\
XGBoost (tuné)      & 0.0001 & 43.33 & 37.59 & $\approx 0$ \\
\bottomrule
\end{tabular}
\end{center}

\smallskip

\begin{block}{Analyse critique}
\small
\begin{itemize}
  \item \textbf{Lasso} : 16/16 coefficients $= 0$ $\rightarrow$ prédit la moyenne (224.3 mg/dL)
  \item \textbf{Data leakage} sur \texttt{hypercholesterol} : $R^2$ 0.675 artificiel, corrigé à $\approx 0$
  \item Aucun modèle retenu : différences statistiquement négligeables
\end{itemize}
\end{block}

\end{frame}

% ─────────────────────────────────────────────────────────────
\begin{frame}{Classification ---~Résultats}

\small\textbf{Protocole :} split stratifié 80/20, K-Fold stratifié ($k=5$), \texttt{class\_weight='balanced'} vs SMOTE

\smallskip

\begin{center}
\small
\begin{tabular}{lrrrrr}
\toprule
Modèle & Accuracy & Précision & Rappel & F1 & AUC \\
\midrule
Naïf (tout à 0) & 0.752 & n/a   & 0.000 & 0.000 & 0.500 \\
Logistic Reg.   & 0.498 & 0.249 & 0.509 & 0.335 & 0.499 \\
Random Forest   & 0.568 & 0.245 & 0.355 & 0.290 & 0.494 \\
XGBoost         & 0.525 & 0.246 & 0.441 & 0.316 & 0.499 \\
\bottomrule
\end{tabular}
\end{center}

\smallskip

\normalsize
\begin{block}{Analyse critique}
\small
\begin{itemize}
  \item \textbf{AUC $\approx$ 0.5} pour tous = équivalent à un tirage aléatoire
  \item SMOTE et \texttt{class\_weight} ne peuvent pas extraire un signal absent
  \item Cohérent avec la régression : \texttt{disease\_risk} généré indépendamment
\end{itemize}
\end{block}

\end{frame}

% ══════════════════════════════════════════════════════════════
\section{Clustering et interprétabilité}
% ══════════════════════════════════════════════════════════════

\begin{frame}{Clustering ---~K-Means et CAH}

\small
\begin{columns}[T]
\begin{column}{0.48\textwidth}
\textbf{K-Means ($k=3$)}

\smallskip
\begin{tabular}{llr}
\toprule
Cluster & Profil & IMC \\
\midrule
0 & Obèse hypertendu & 35.0 \\
1 & Surpoids normotendu & 29.1 \\
2 & Normal hypertendu & 24.0 \\
\bottomrule
\end{tabular}

\smallskip
Partition : \textbf{IMC} + \textbf{pression systolique}
\end{column}

\begin{column}{0.48\textwidth}
\textbf{CAH Ward ($n=5\,000$)}

\smallskip
\begin{tabular}{llr}
\toprule
Cluster & Profil & Fumeurs \\
\midrule
0 & Non-fumeurs, HTA & 2\% \\
1 & Fumeurs modérés & 18\% \\
2 & Fumeurs intenses, HTA & \textbf{100\%} \\
\bottomrule
\end{tabular}

\smallskip
Isole un groupe de \textbf{fumeurs intenses}
\end{column}
\end{columns}

\medskip

\begin{tabular}{lrrl}
\toprule
Métrique & K-Means & CAH & \\
\midrule
Silhouette & 0.103 & 0.105 & Faible ---~ données sans structure \\
ACP 2D     & \multicolumn{2}{c}{24.3\% de variance} & 15 dims quasi-indépendantes \\
\bottomrule
\end{tabular}

\end{frame}

% ─────────────────────────────────────────────────────────────
\begin{frame}{Interprétabilité ---~Convergence des méthodes}

\textbf{5 méthodes :} coefficients LR, Ridge, feature importance RF \& XGBoost, SHAP

\medskip

\begin{center}
\small
\begin{tabular}{lcccccr}
\toprule
Variable & LR & RF & XGB & SHAP & Ridge & Score \\
\midrule
\texttt{systolic\_bp}    & \checkmark & \checkmark & \checkmark & \checkmark & \checkmark & \textbf{5/5} \\
\texttt{daily\_steps}    & \checkmark & \checkmark & $\cdot$    & \checkmark & $\cdot$    & 3/5 \\
\texttt{family\_history} & \checkmark & $\cdot$    & \checkmark & $\cdot$    & \checkmark & 3/5 \\
\texttt{bmi}             & $\cdot$    & \checkmark & $\cdot$    & \checkmark & $\cdot$    & 2/5 \\
\bottomrule
\end{tabular}
\end{center}

\medskip

\normalsize
\begin{block}{Interprétation}
\begin{itemize}
  \item Les 5 méthodes \textbf{ne convergent pas} $\rightarrow$ elles mesurent du bruit aléatoire
  \item \texttt{systolic\_bp} partout : grande variance, pas signal causal
  \item \textbf{La divergence est elle-même un résultat analytique}
\end{itemize}
\end{block}

\end{frame}

% ══════════════════════════════════════════════════════════════
\section{Bilan et perspectives}
% ══════════════════════════════════════════════════════════════

\begin{frame}{Bilan ---~Résultat central}

\begin{block}{\alert{Limite fondamentale du dataset synthétique}}
Les variables ont été générées \textbf{indépendamment}, sans relation causale encodée.\\
Aucune technique ne peut extraire un signal qui n'existe pas.
\end{block}

\bigskip

\begin{columns}[T]
\begin{column}{0.48\textwidth}
\textbf{Ce que ça prouve :}
\begin{itemize}
  \item Corrélations $< 0.01$ en EDA prédisaient ce résultat
  \item Lasso : 16/16 coefficients $= 0$
  \item $R^2 \approx 0$, AUC $\approx 0.5$ partout
  \item Interprétabilité divergente
\end{itemize}
\end{column}

\begin{column}{0.48\textwidth}
\textbf{Ce n'est pas un échec :}
\begin{itemize}
  \item Data leakage détecté et corrigé
  \item Pipeline rigoureux et complet
  \item Diagnostic critique mené à chaque étape
  \item Directement transposable à des données réelles
\end{itemize}
\end{column}
\end{columns}

\bigskip

\textcolor{bleuaccent}{\textbf{Savoir diagnostiquer pourquoi un modèle échoue est aussi précieux que d'en construire un qui réussit.}}

\end{frame}

% ─────────────────────────────────────────────────────────────
\begin{frame}{Perspectives}

\begin{enumerate}
  \item \textbf{Changer de dataset} ---~cohortes Framingham ou NHANES : vraies relations biologiques, même pipeline
  \item \textbf{Améliorer la classification} ---~stacking, réseaux de neurones (Deep Learning, M2)
  \item \textbf{Déploiement} ---~API REST (FastAPI) avec prédictions + valeurs SHAP associées
  \item \textbf{Interprétabilité avancée} ---~Partial Dependence Plots (PDP), ICE
  \item \textbf{Data drift} ---~suivi de l'AUC sur nouvelles cohortes, réentraînement périodique
  \item \textbf{Master 2} ---~causalité, fairness algorithmique, apprentissage fédéré
\end{enumerate}

\end{frame}

% ── Diapo finale ─────────────────────────────────────────────
\begin{frame}

\vspace{2cm}
\begin{center}
{\Large\textcolor{bleutitre}{\textbf{Merci pour votre attention}}}

\bigskip\bigskip

{\large Questions ?}

\bigskip\bigskip\bigskip

\textcolor{gray}{\small Ardacham Mahamat Teguene ---~M1 MIASHS ---~Juin 2026}
\end{center}

\end{frame}

\end{document}
